\subsection{\texorpdfstring{PRIMITIVE ELEMENTS OF \(\mathfrak{sl}(2,k)\) -MODULES}{PRIMITIVE ELEMENTS OF \textbackslash mathfrak\{sl\}(2,k) -MODULES}}

Let E be an \(\mathfrak{sl}(2,k)\) -module. If \(A \in \mathfrak{sl}(2,k)\) and \(x \in E\) , we shall often write Ax instead of \(A_{E}x\) . Let \(\lambda \in k\) . If Hx = \(\lambda x\) we say, by abuse of language, that x is an element of E of weight \(\lambda\) , or that \(\lambda\) is the weight of x. If E is finite dimensional, \(H_{E}\) is semi-simple, so the set of elements of weight \(\lambda\) is the primary subspace of E relative to \(H_{E}\) and \(\lambda\) (cf.~Chap. VII, §1, no. 1).

Lemma 3. If \(x\) is an element of weight \(\lambda\) , then \(X_{+}x\) is an element of weight \(\lambda + 2\) and \(X_{-}x\) is an element of weight \(\lambda - 2\) .

Indeed, \(HX_{+}x = [H,X_{+}]x + X_{+}Hx = 2X_{+}x + X_{+}\lambda x = (\lambda +2)X_{+}x\) , and similarly \(HX_{-}x = (\lambda -2)X_{-}x\) (cf.~also Chap. VII, §1, no. 3, Prop. 10 (ii)).

DEFINITION 1. Let E be an \(\mathfrak{sl}(2,k)\) -module. An element of E is said to be primitive if it is a non-zero eigenvector of \(H_{E}\) and belongs to the kernel of \(X_{+E}\) .

A non-zero element e of E is primitive if and only if ke is stable under the operation of \(kH + kX_{+}\) ; this follows for example from Lemma 3.

Examples The element \(X_{+}\) is primitive of weight 2 for the adjoint representation of \(\mathfrak{sl}(2,k)\) . The element \((1,0)\) of \(k^{2}\) is primitive of weight 1 for the identity representation of \(\mathfrak{sl}(2,k)\) on \(k^{2}\) .

Lemma 4. Let \(\mathbf{E}\) be a non-zero finite dimensional \(\mathfrak{sl}(2,k)\) -module. Then \(\mathbf{E}\) has primitive elements.

Since \(X_{+}\) is a nilpotent element of \(\mathfrak{sl}(2,k)\) , \(X_{+E}\) is nilpotent. Assume that \(X_{+E}^{m-1} \neq 0\) and \(X_{+E}^{m} = 0\) . By Lemma 2,

\[
m (H _ {\mathrm{E}} - m + 1) X _ {+ \mathrm{E}} ^ {m - 1} = [ X _ {- \mathrm{E}}, X _ {+ \mathrm{E}} ^ {m} ] = 0,
\]

and hence the elements of \(X_{+}^{m - 1}(\mathrm{E}) - \{0\}\) are primitive.

PROPOSITION 1. Let E be an \(\mathfrak{sl}(2,k)\) -module, and \(e\) a primitive element of E of weight \(\lambda\) . Put \(e_n = \frac{(-1)^n}{n} X_-^n e\) for \(n \geq 0\) , and \(e_{-1} = 0\) . Then

\[
\left\{ \begin{array}{l l} H e _ {n} & = (\lambda - 2 n) e _ {n} \\ X _ {-} e _ {n} & = - (n + 1) e _ {n + 1} \\ X _ {+} e _ {n} & = (\lambda - n + 1) e _ {n - 1}. \end{array} \right.\tag{2}
\]

The first formula follows from Lemma 3, and the second from the definition of the \(e_n\) . We prove the third by induction on \(n\) . It is satisfied for \(n = 0\) since \(e_{-1} = 0\) . If \(n > 0\) ,

\[
\begin{array}{r l} & n X _ {+} e _ {n} = - X _ {+} X _ {-} e _ {n - 1} = - [ X _ {+}, X _ {-} ] e _ {n - 1} - X _ {-} X _ {+} e _ {n - 1} \\ & \qquad = H e _ {n - 1} - X _ {-} (\lambda - n + 2) e _ {n - 2} \\ & \qquad = (\lambda - 2 n + 2 + (n - 1) (\lambda - n + 2)) e _ {n - 1} \\ & \qquad = n (\lambda - n + 1) e _ {n - 1}. \end{array}
\]

COROLLARY. The submodule of \(\mathbf{E}\) generated by \(e\) is the vector subspace generated by the \(e_n\) .

This follows from the formulas (2).

The integers \(n \geq 0\) such that \(e_{n} \neq 0\) constitute an interval in N, and the corresponding elements \(e_{n}\) form a basis over k of the submodule generated by e (indeed, they are linearly independent because they are non-zero elements of distinct weights). This basis will be said to be associated to the primitive element e.

PROPOSITION 2. If the submodule V of E generated by the primitive element e is finite dimensional, then:

\begin{enumerate}
\def\labelenumi{(\roman{enumi})}
\item
  the weight \(\lambda\) of e is integral and equal to \(\dim V - 1\) ;
\item
  \((e_0, e_1, \ldots, e_\lambda)\) is a basis of V, and \(e_n = 0\) for \(n > \lambda\) ;
\item
  the eigenvalues of \(H_{\mathrm{V}}\) are \(\lambda, \lambda - 2, \lambda - 4, \ldots, -\lambda\) ; they are all of multiplicity 1;
\item
  every primitive element of V is proportional to e;
\item
  the commutant of the module \(\mathrm{V}\) is reduced to the scalars; in particular, \(\mathrm{V}\) is absolutely simple.
\end{enumerate}

Let \(m\) be the largest integer such that \(e_m \neq 0\) . Then \(0 = X_{+}e_{m+1} = (\lambda - m)e_m\) , so \(\lambda = m\) ; since \((e_0, e_1, \ldots, e_m)\) is a basis of \(V\) , this proves (i) and (ii). Assertion (iii) follows from the equality \(He_n = (\lambda - 2n)e_n\) . We have \(X_{+}e_n \neq 0\) for \(1 \leq n \leq m\) , hence (iv). Let \(c\) be an element of the commutant of the module \(V\) . Then \(Hc(e) = cH(e) = \lambda c(e)\) , so there exists \(\mu \in k\) such that \(c(e) = \mu e\) ; then

\[
c X _ {-} ^ {q} e = X _ {-} ^ {q} c e = \mu X _ {-} ^ {q} e
\]

for all \(q \geq 0\) , so \(c = \mu.1\) , proving (v).

COROLLARY. Let E be a finite dimensional \(\mathfrak{sl}(2,k)\) -module.

\begin{enumerate}
\def\labelenumi{(\roman{enumi})}
\item
  The endomorphism \(H_{E}\) is diagonalizable and its eigenvalues are rational integers.
\item
  For any \(p \in Z\) , let \(E_{p}\) be the eigenspace of \(H_{E}\) corresponding to the eigenvalue p.~Let i be an integer \(\geq 0\) . The map \(X_{-E}^{i}|E_{p}:E_{p} \to E_{p-2i}\) is injective for \(i \leq p\) , bijective for i = p, and surjective for \(i \geq p\) . The map \(X_{+E}^{i}|E_{-p}:E_{-p} \to E_{-p+2i}\) is injective for \(i \leq p\) , bijective for i = p, and surjective for \(i \geq p\) .
\item
  The length of \(\mathrm{E}\) is equal to \(\dim \operatorname{Ker} X_{+\mathrm{E}}\) and to \(\dim \operatorname{Ker} X_{-\mathrm{E}}\) .
\item
  Let \(E'\) (resp. \(E''\) ) be the sum of the \(E_{p}\) for p even (resp. odd). Then \(E'\) (resp. \(E''\) ) is the sum of the simple submodules of E of odd (resp. even) dimension; and \(E = E' \oplus E''\) . The length of \(E'\) is \(\dim E_{0}\) , and that of \(E''\) is \(\dim E_{1}\) .
\end{enumerate}

\[
\text {(v)} \operatorname{Ker} X _ {+ \mathrm{E}} \cap \operatorname{Im} X _ {+ \mathrm{E}} \subset \sum_ {p > 0} \mathrm{E} _ {p} \text {and} \operatorname{Ker} X _ {- \mathrm{E}} \cap \operatorname{Im} X _ {- \mathrm{E}} \subset \sum_ {p <   0} \mathrm{E} _ {p}.
\]

If E is simple, E is generated by a primitive element (Lemma 4), and it suffices to apply Propositions 1 and 2. The general case follows since every finite dimensional \(\mathfrak{sl}(2,k)\) -module is semi-simple.

